% GENERATED FILE. Edit canonical lessons, then run npm run generate:books.
% canonical-source-hash: fnv1a64:6fab28f025b7827a
% canonical-lessons: ZH-C244-taifeng, ZH-C244-zhenyu, ZH-C244-shuobuding, ZH-C244-wanyi, ZH-C244-jixu

\chapter{A Storm Coming}
\label{ch:zh-a-storm-coming}
\hlchaptermodality{244}

\begin{chapteropening}
\textbf{By the end of this chapter:} \emph{I can talk about plans when a typhoon or showers may come: who knows, in case, and carrying on.}

Complete the last lesson of chapter 244: I can talk about plans when a typhoon or showers may come: who knows, in case, and carrying on.
\end{chapteropening}

\section[táifēng]{táifēng — typhoon}

\label{lesson:ZH-C244-taifeng}



\begin{quote}
Before the new one: say the Mandarin for to freeze over, then the Mandarin for temperature.
\end{quote}

\subsection*{You'll want to know: táifēng}
\emph{táifēng} — \textquotedblleft{}typhoon\textquotedblright{}.

\begin{quote}
\emph{hòutiān huì yǒu táifēng} — There'll be a typhoon the day after tomorrow.
\end{quote}

Its characters come later; for now it is a word for the ear.

\begin{cousinweb}
\emph{táifēng} is built from \emph{tái} and \zh{风} 'wind'. The second is inherited from Middle Chinese \emph{pjuwng} (Guangyun, in Baxter's transcription); no dictionary we use records an older form of the first in this sense, so its history is not settled.
\end{cousinweb}

\begin{grammarlens}[title={What you've built}]
One more word for talking about plans.
\end{grammarlens}

\subsection*{Your turn}
\begin{itemize}
\raggedright
  \item \emph{Say it:} \emph{táifēng}
  \item \emph{Say it:} \emph{táifēng}, once more
  \item \emph{Say it:} \emph{hòutiān huì yǒu táifēng} — There'll be a typhoon the day after tomorrow.
  \item \emph{Recall:} say the Mandarin for to freeze over, then the Mandarin for temperature, then say \emph{táifēng} again
\end{itemize}

\begin{tcolorbox}[breakable,colback=teal!4,colframe=teal!35!black,title={Before you move on}]
What does táifēng mean? (\textquotedblleft{}Typhoon\textquotedblright{}.) Say it once more.
\end{tcolorbox}
\section[zhènyǔ]{zhènyǔ — a rain shower}

\label{lesson:ZH-C244-zhenyu}



\begin{quote}
Before the new one: say the Mandarin for temperature, then the Mandarin for a typhoon.
\end{quote}

\subsection*{You'll want to know: zhènyǔ}
\emph{zhènyǔ} — \textquotedblleft{}a rain shower\textquotedblright{}.

\begin{quote}
\emph{míngtiān xiàwǔ yǒu zhènyǔ} — There'll be showers tomorrow afternoon.
\end{quote}

Its characters come later; for now it is a word for the ear.

\begin{cousinweb}
\emph{zhènyǔ} is built from \emph{zhèn} 'a spell' and \zh{雨} 'rain', both inherited from Middle Chinese, \emph{drinH} and \emph{hjuX} (Guangyun, in Baxter's transcription); the dictionary lists the first only as the everyday form of another character, and the sense 'a spell' is not given there.
\end{cousinweb}

\begin{grammarlens}[title={What you've built}]
One more word for talking about plans.
\end{grammarlens}

\subsection*{Your turn}
\begin{itemize}
\raggedright
  \item \emph{Say it:} \emph{zhènyǔ}
  \item \emph{Say it:} \emph{zhènyǔ}, once more
  \item \emph{Say it:} \emph{míngtiān xiàwǔ yǒu zhènyǔ} — There'll be showers tomorrow afternoon.
  \item \emph{Recall:} say the Mandarin for temperature, then the Mandarin for a typhoon, then say \emph{zhènyǔ} again
\end{itemize}

\begin{tcolorbox}[breakable,colback=teal!4,colframe=teal!35!black,title={Before you move on}]
What does zhènyǔ mean? (\textquotedblleft{}A rain shower\textquotedblright{}.) Say it once more.
\end{tcolorbox}
\section[shuōbudìng]{shuōbudìng — perhaps; who knows}

\label{lesson:ZH-C244-shuobuding}



\begin{quote}
Before the new one: say the Mandarin for a typhoon, then the Mandarin for a rain shower.
\end{quote}

\subsection*{You'll want to know: shuōbudìng}
\emph{shuōbudìng} — \textquotedblleft{}perhaps; who knows\textquotedblright{}.

\begin{quote}
\emph{shuōbudìng míngtiān huì xiàyǔ} — Who knows, it might rain tomorrow.
\end{quote}

Its characters come later; for now it is a word for the ear.

Tone: \zh{不} is light here, bu, with no tone of its own.

\begin{cousinweb}
\emph{shuōbudìng} is built from \zh{说} 'to say', \zh{不} 'not' and \emph{dìng} 'fixed', all inherited from Middle Chinese, \emph{sywet}, \emph{pjut} and \emph{dengH} (Guangyun, in Baxter's transcription), which glosses the last 'settled'.
\end{cousinweb}

\begin{grammarlens}[title={What you've built}]
One more word for talking about plans.
\end{grammarlens}

\subsection*{Your turn}
\begin{itemize}
\raggedright
  \item \emph{Say it:} \emph{shuōbudìng}
  \item \emph{Say it:} \emph{shuōbudìng}, once more
  \item \emph{Say it:} \emph{shuōbudìng míngtiān huì xiàyǔ} — Who knows, it might rain tomorrow.
  \item \emph{Recall:} say the Mandarin for a typhoon, then the Mandarin for a rain shower, then say \emph{shuōbudìng} again
\end{itemize}

\begin{tcolorbox}[breakable,colback=teal!4,colframe=teal!35!black,title={Before you move on}]
What does shuōbudìng mean? (\textquotedblleft{}Perhaps; who knows\textquotedblright{}.) Say it once more.
\end{tcolorbox}
\section[wànyī]{wànyī — in case; if by any chance}

\label{lesson:ZH-C244-wanyi}



\begin{quote}
Before the new one: say the Mandarin for a rain shower, then the Mandarin for who knows.
\end{quote}

\subsection*{You'll want to know: wànyī}
\emph{wànyī} — \textquotedblleft{}in case; if by any chance\textquotedblright{}.

\begin{quote}
\emph{wànyī xiàyǔ zěnme bàn} — What if it rains?
\end{quote}

Its characters come later; for now it is a word for the ear.

Tone: \zh{一} keeps its first tone here, yī, because it ends the word.

\begin{cousinweb}
\emph{wànyī} is built from \emph{wàn} 'ten thousand' and \zh{一} 'one', both inherited from Middle Chinese, \emph{mjonH} and \emph{'jit} (Guangyun, in Baxter's transcription); it is usually explained as one chance in ten thousand.
\end{cousinweb}

\begin{grammarlens}[title={What you've built}]
One more word for talking about plans.
\end{grammarlens}

\subsection*{Your turn}
\begin{itemize}
\raggedright
  \item \emph{Say it:} \emph{wànyī}
  \item \emph{Say it:} \emph{wànyī}, once more
  \item \emph{Say it:} \emph{wànyī xiàyǔ zěnme bàn} — What if it rains?
  \item \emph{Recall:} say the Mandarin for a rain shower, then the Mandarin for who knows, then say \emph{wànyī} again
\end{itemize}

\begin{tcolorbox}[breakable,colback=teal!4,colframe=teal!35!black,title={Before you move on}]
What does wànyī mean? (\textquotedblleft{}In case; if by any chance\textquotedblright{}.) Say it once more.
\end{tcolorbox}
\section[jìxù]{jìxù — to carry on; to continue}

\label{lesson:ZH-C244-jixu}



\begin{quote}
Before the new one: say the Mandarin for who knows, then the Mandarin for in case.
\end{quote}

\subsection*{You'll want to know: jìxù}
\emph{jìxù} — \textquotedblleft{}to carry on; to continue\textquotedblright{}.

\begin{quote}
\emph{yǔ tíng le wǒmen jìxù zǒu} — The rain's stopped; let's carry on.
\end{quote}

Its characters come later; for now it is a word for the ear.

\begin{cousinweb}
\emph{jìxù} is built from \emph{jì} 'to continue' and \emph{xù} 'to join on'. The second is inherited from Middle Chinese \emph{zjowk} (Guangyun, in Baxter's transcription), glossed 'to continue' and 'to join'; the dictionary glosses the first, \emph{kejH} there, 'to carry on'.
\end{cousinweb}

\begin{grammarlens}[title={What you've built}]
One more word for talking about plans.
\end{grammarlens}

\subsection*{Your turn}
\begin{itemize}
\raggedright
  \item \emph{Say it:} \emph{jìxù}
  \item \emph{Say it:} \emph{jìxù}, once more
  \item \emph{Say it:} \emph{yǔ tíng le wǒmen jìxù zǒu} — The rain's stopped; let's carry on.
  \item \emph{Recall:} say the Mandarin for who knows, then the Mandarin for in case, then say \emph{jìxù} again
\end{itemize}

\begin{tcolorbox}[breakable,colback=teal!4,colframe=teal!35!black,title={Before you move on}]
What does jìxù mean? (\textquotedblleft{}To carry on; to continue\textquotedblright{}.) Say it once more.
\end{tcolorbox}
